\documentclass[12pt,examen]{repaso_layout}
\campoformativo{Saberes y Pensamiento Científico}
\grado{2}
\nivel{Secundaria}
\materia{Matemáticas}
\unidad{1}
\title{Extraordinario}
\aprendizajes{
      \item Resuelve problemas que impliquen la suma, resta, la multiplicación y la división de números enteros, aplicando las reglas correspondientes.
      \item Identifica y ubica números negativos en una recta numérica, comparando su magnitud.
      \item Aplica las propiedades de las potencias a números negativos en la resolución de problemas.
      \item Resuelve problemas que involucren las leyes de los exponentes, y expresa números en notación científica.
      \item Resuelve problemas de contexto científico y tecnológico utilizando la notación científica.
      \item Identifica y ubica puntos en el plano cartesiano, y comprende la estructura de los cuadrantes.
      \item Calcula la pendiente de una recta y comprende su significado en diferentes contextos.
      \item Resuelve problemas que involucren el uso de porcentajes, mediante la ``regla de tres''.
}
\author{Melchor Pinto, JC}
\begin{document}
    HJola, qué hace?
\begin{questions}
\section{Sección 1}
\subsection{Subsección 1}
\question[10]{Hola}
\subsection{Subsección 2}
\question[10]{Hola}
\section{Sección 2}
\question[10]{Hola}
\question[10]{Hola}
\subsection{Subsección 1}
\question[10]{Hola}
\subsection{Subsección 2}
\question[10]{Hola}
\subsection{Subsección 3}
\question[10]{Hola}
\question[10]{Hola}
\subsection{Subsección 4}
\question[10]{Hola}
\section{Sección 3} 
\question[10]{Hola}
\question[10]{Hola}
\subsection{Subsección 1}
\question[10]{Hola}
\subsection{Subsección 2}
\question[10]{Hola}
\question[10]{Hola}
\question[10]{Hola}
\subsection{Subsección 3}
\question[10]{Hola}
\question[10]{Hola}
\subsection{Subsección 4}
\question[10]{Hola}
\question[10]{Hola}
\section{Sección 4}
\question[10]{Hola}
\question[10]{Hola}
\subsection{Subsección 1}
\question[10]{Hola}
\question[10]{Hola}
\question[10]{Hola}
\question[10]{Hola}
\subsection{Subsección 2}
\question[10]{Hola}
\section{Sección 5}
\question[10]{Hola}
\question[10]{Hola}
\subsection{Subsección 1}
\question[10]{Hola}
\subsection{Subsección 2}
\question[10]{Hola}
\end{questions}
\end{document}